%======================================================================
\chapter{Introduction}\label{ch:intro}
%======================================================================

\section{What this volume is about}

Volumes~1 and~2 dealt with one object, the pseudo-arc. The rest of the classification of
homogeneous plane continua needs a stock of general facts about continua, and this volume
collects them. Most of them are classical and can be found, in some form, in textbooks on
continuum theory; here they are proved in the form in which the formal proof uses them, and
the proofs follow the formal ones.

The volume has seven chapters, which can be read almost independently.

\begin{itemize}
\item \textbf{Chapter~\ref{ch:bumping}: Separated sets and boundary bumping.} Separated sets, the
union of a continuum with one piece of a separation of its complement, nested intersections of
continua, components as intersections of clopen sets, and the \emph{boundary bumping theorem}:
in a continuum $M$, the component of a point $p$ in a proper closed subset $F$ reaches the
boundary of $F$.
\item \textbf{Chapter~\ref{ch:irreducible}: Irreducible continua.} A continuum containing two
points (or meeting two compact sets) contains a subcontinuum that is minimal with this
property; a nondegenerate indecomposable continuum is irreducible between two of its points.
\item \textbf{Chapter~\ref{ch:noncut}: Non-cut points.} R.~L.~Moore's theorem: a nondegenerate
compact connected $T_1$ space has at least two points whose removal leaves it connected, and no
proper connected subset contains all such points.
\item \textbf{Chapter~\ref{ch:triods}: Triods and Sorgenfrey's theorem.} Triods, weak triods and
unicoherence; an indecomposable continuum is not a triod; and Sorgenfrey's theorem: a
nondegenerate unicoherent continuum which is not a triod is irreducible between two of its
points.
\item \textbf{Chapter~\ref{ch:snake}: Snake-like and circle-like continua.} Subcontinua of
snake-like continua are snake-like; snake-like continua are unicoherent and contain no triod;
the $\varepsilon$-map characterization of snake-likeness and its consequence that snake-likeness
is a topological property; circular chains, the circle, and proper subcontinua of circle-like
continua.
\item \textbf{Chapter~\ref{ch:decomp}: Upper semicontinuous decompositions.} Decompositions of a
compact metric space into closed sets, the decomposition space and its metrizability, monotone
decompositions, homeomorphisms permuting the elements, and the theorem of Kuratowski and Fort
that an upper semicontinuous decomposition is continuous at the elements of a residual set.
\item \textbf{Chapter~\ref{ch:effros}: The homeomorphism group and the Effros theorem.} The group
of self-homeomorphisms of a compact metric space is a Polish group; a Polish group acting
transitively on a Baire space acts micro-transitively (the open mapping principle, following
J.~van Mill's elementary proof); and Ungar's corollary for homogeneous compact metric spaces:
close points are moved one to the other by homeomorphisms close to the identity.
\end{itemize}

The main results, in the form used later, are:

\begin{theorem}[Boundary bumping; Theorem~\ref{thm:bump-closed}]
Let $M$ be a continuum, $F\subsetneq M$ a closed proper subset and $p\in F$. Some subcontinuum of
$F$ containing $p$ meets $\cl{M\setminus F}$.
\end{theorem}

\begin{theorem}[Moore; Theorem~\ref{thm:two-noncut}]
A compact connected $T_1$ space with at least two points has two distinct non-cut points.
\end{theorem}

\begin{theorem}[Sorgenfrey; Theorem~\ref{thm:sorgenfrey}]
A nondegenerate unicoherent continuum that is not a triod is irreducible between two of its
points.
\end{theorem}

\begin{theorem}[Theorems~\ref{thm:snake-unicoherent} and~\ref{thm:snake-atriodic}]
A snake-like continuum is hereditarily unicoherent and atriodic.
\end{theorem}

\begin{theorem}[Theorems~\ref{thm:quotient-metrizable} and~\ref{thm:residual}]
The decomposition space of an upper semicontinuous decomposition of a compact metric space into
closed sets is a compact metrizable space, and the decomposition is continuous at the elements of
a residual subset of it.
\end{theorem}

\begin{theorem}[Effros, Ungar; Theorem~\ref{thm:ungar}]
Let $Y$ be a homogeneous compact metric space. For every $\varepsilon>0$ there is $\delta>0$ such
that whenever $d(x,y)<\delta$ there is a homeomorphism $h\colon Y\to Y$ with $h(x)=y$ and
$d(h(z),z)<\varepsilon$ for all $z\in Y$.
\end{theorem}

\section{Where these results are used}

The results of this volume are used throughout Volumes~4--8. The most important uses are:
boundary bumping, in almost every argument about plane continua; irreducibility and Sorgenfrey's
theorem, in the study of snake-like and circle-like continua (Fugate's and Ingram's theorems);
non-cut points, for the end elements of Bing's decomposition of an atriodic continuum; the
fact that indecomposable continua are not triods, together with Moore's theorem on triods in the
plane (a theorem of plane topology), in Hagopian's proof that a homogeneous indecomposable plane continuum is
atriodic; snake-likeness and circle-likeness, in the indecomposable case and for circles of
pseudo-arcs; upper semicontinuous decompositions, for Jones's aposyndetic decomposition of a
decomposable homogeneous continuum (Volume~7); and the Effros theorem, wherever homogeneity is
used quantitatively.

\section{Where the formal proof differs from the literature}\label{sec:intro-diff}

The notes follow the formal proof. It departs from the usual sources in the following places.

\begin{enumerate}
\item \textbf{Non-cut points without a metric.} Moore's theorem is proved, as in Nadler's
\emph{Continuum Theory} (6.3--6.7), for compact connected $T_1$ spaces; the formal statements drop a
hypothesis of nondegeneracy which is implied by the existence of a cut point
(Chapter~\ref{ch:noncut}).
\item \textbf{A false lemma in the proof of Sorgenfrey's theorem.} The lemma that removing a
subcontinuum $H$ from the set $U=X\setminus Z$ of an essential sum leaves a connected set
(Nadler~11.28) is false when $H$ contains $U$: the empty set is not connected. The proof assumes,
without saying so, that $U\setminus H$ is nonempty, and the formal proof adds this hypothesis
(Remark~\ref{rem:1128-false}).
\item \textbf{Sorgenfrey's induction simplified.} In the main step of Sorgenfrey's theorem
(Theorem~\ref{thm:irr-point}) the formal proof absorbs, at stage $n$, the $n$-th point of a dense
sequence, instead of the first point not yet absorbed; only the properties listed in
Lemma~\ref{lem:sorg-step} are used afterwards. Two hypotheses of Nadler's 11.31 and 11.32 are not
needed and are dropped.
\item \textbf{Triods from a core.} In place of Sorgenfrey's theorem that a weak triod contains a
triod, which the formal proof does not use (and does not prove), the later volumes use a simple
``core'' construction (Lemma~\ref{lem:core}).
\item \textbf{Snake-likeness in a compact space.} Snake-likeness is defined with chains of open
sets of the ambient space. Its invariance under homeomorphisms (Theorem~\ref{thm:snake-homeo}) is
proved through $\varepsilon$-maps and the Tietze extension theorem, and needs the target space to
be compact.
\item \textbf{The Effros theorem.} Van Mill's Proposition~2.2 (two analytic fat sets meet) needs
the space to be nonempty; the formal proof proves a relative form (Theorem~\ref{thm:prop22}),
which is what Proposition~3.5 uses. Van Mill's Lemma~3.4 and Proposition~3.5 are proved as
printed, with closures.
\end{enumerate}

\section{Prerequisites and conventions}

The reader should know the basic metric topology of \vone{Chapter~1} (compactness,
connectedness, continua, indecomposability) and, for Chapter~\ref{ch:snake}, the definition of a
chain (\vone{Definition~2.1}). Chapter~\ref{ch:irreducible} uses one result of Volume~2: in an
indecomposable continuum no composant is the whole continuum (\vtwo{Theorem~1.4}).

\begin{itemize}
\item $X$, $S$, $Y$ denote metric spaces unless stated otherwise, $d$ their metrics. $\cl A$ is the
closure of $A$ and $\interior A$ its interior; $B(x,r)$ and $\bar B(x,r)$ are the open and closed
balls.
\item A \emph{continuum} is a nonempty compact connected set (\vone{Definition~1.2}). In the
formal proof this is \lean{HPC.Fold.Q78.IsContinuum}. Most chapters work in a \emph{compact}
metric space $S$; then a set is a continuum if and only if it is nonempty, closed and connected.
\item \emph{Preconnected} means connected or empty, as in \vone{Section~1.3}; the criterion
\vone{(K2)} for closed sets is used constantly: a closed set $C$ is preconnected if and only if
it is not the union of two disjoint nonempty closed sets.
\item A set is \emph{clopen} if it is both closed and open. A subset $W$ of a set $F$ is
\emph{relatively clopen in $F$} if it is clopen in the subspace $F$.
\item $\N=\set{0,1,2,\dots}$ and $[n]=\set{0,\dots,n-1}$.
\end{itemize}

\paragraph{Formal counterparts.} As in Volumes~1 and~2, many results are followed by a line
``Formal counterpart:'' naming the corresponding declaration of the Lean development (all names
begin with \lean{HPC.}). These lines may be skipped. Appendix~\ref{app:formal} collects them.

\paragraph{Exercises.} Each chapter ends with exercises. None of them is used later.
